\chapter{Canonical Evidence Encoding}\label{ch05:encoding}

\section{Why digests are load-bearing}\label{ch05:loadbearing}

Every mechanism described in this monograph eventually reduces a question
about meaning to a question about bytes. A semantic reference binds a concept
to the digest of its defining text (Chapter 3). A registry record carries a
digest bound to the ledger bytes it was parsed from (Chapter 6). A
compatibility receipt binds compiler bytes, worker bytes, and a result digest
(Chapter 7), and the replay requirement compares fresh digests against
recorded ones before any structural decision can pass (Chapter 9). In each
case the operative comparison is hash equality of encoded values: two
artifacts are treated as the same evidence exactly when their canonical
encodings hash identically.

This reduction is what makes the kernel auditable---byte equality is a
question any party can settle without trusting the kernel---but it places the
entire weight of the system on one function: the encoder that turns in-memory
values into the bytes that get hashed. If that function is nondeterministic,
platform-dependent, or many-to-one in the wrong places, then every downstream
comparison is silently invalid, and no amount of care in the layers above
repairs it. The failure is maximally quiet: hashes still compute, receipts
still verify against themselves, tests still pass on the machine that wrote
them, and the divergence surfaces---if it surfaces at all---as an
irreproducible mismatch between two honest parties, months later, with the
original context gone.

\begin{definition}[Canonical evidence encoding]\label{ch05:def:encoding}
The \emph{canonical evidence encoding} is the kernel's single function from
admissible in-memory values to byte strings, with the contract that (i) equal
admissible values encode to equal bytes on every platform and every run
(determinism), (ii) values the kernel treats as semantically distinct encode
to distinct bytes (injectivity on the admitted domain), and (iii) every value
outside the admissible domain is rejected with a typed error rather than
encoded by best effort (fail-closed totality).
\end{definition}

The three clauses are jointly the definition of ``canonical'' as this
monograph uses the word. Clause (i) rules out dependence on iteration order,
hash seeds, locale, word size, and library version. Clause (ii) rules out
\emph{type puns}: pairs of distinct values that encode identically, so that
evidence about one can be replayed as evidence about the other. Clause (iii)
is the clause most often skipped in practice and the one the kernel is most
insistent about: an encoder that guesses on hard inputs converts every hard
input into a latent comparison bug, whereas an encoder that rejects them
converts the same inputs into visible errors at the moment of encoding, when
the context needed to fix them still exists.

The design follows the same instinct as the reproducible-builds
discipline \cite{reprobuilds}: identity is a function
of bytes, so byte production must be an engineered, specified, testable
process rather than an incidental property of whatever serializer was handy.
And, per the standing distinction of Chapter 4, none of this confers
authority. A canonical digest says \emph{these bytes}; it never says
\emph{these bytes are approved}. The encoding chapter is about making the
first statement exact so that the second statement, if a future reviewed
design ever makes it, has a definite object.

\begin{boundary}[Chapter scope]\label{ch05:bd:scope}
This chapter specifies the v0.1 encoding profile and the rationale for each
rule, with worked encodings. Statements that the kernel ``rejects'' or
``produces'' particular bytes are statements about the v0.1 candidate's
tested behavior on the Gate-4 fixture families (acceptance and rejection
cases for sets, floats, NaN, infinity, decimal tagging, cycles, and depth
bounds are among the enumerated suite contents); they are bounded software
receipts in the sense of Chapter 12, not conformance claims over untested
inputs, and the profile itself remains a Ring-4 research candidate.
\end{boundary}

\section{The threat model for encodings}\label{ch05:threats}

Before the rules, the failures they answer. Encoding failures come in two
dual families, and it is worth having both clearly in view because several
rules exist to block one family without reintroducing the other.

\begin{threat}[Digest instability: same evidence, different bytes]
\label{ch05:threat:instability}
Two encodings of the same value differ---across runs, platforms, library
versions, or processes---so hash comparison reports difference where there is
none. Consequences: replay verification fails spuriously; deduplication
splits; honest parties cannot reproduce each other's digests; and, worst,
operators learn to treat digest mismatches as noise, at which point the
mismatch that matters is waved through. Classic causes: unordered container
iteration, hash-seed-dependent ordering, float formatting drift, locale-
dependent case mapping, newline rewriting in transit, Unicode normalization
applied by one layer and not another.
\end{threat}

\begin{threat}[Type puns: different evidence, same bytes]
\label{ch05:threat:puns}
Two values the system treats as semantically distinct encode to the same
bytes, so a digest computed over one validates the other. Consequences: an
evidence record can be constructed to collide with a different record;
quantities can be silently reinterpreted (an exact decimal read back as a
string, a tuple as a list, an enumeration member as its underlying value);
and any downstream logic keyed on digests inherits the confusion. Classic
causes: encoders that erase type distinctions on output (most JSON encoders
erase tuple/list, and many erase decimal/float/string distinctions), and
container coercions applied before encoding.
\end{threat}

The two threats pull in opposite directions. Aggressive normalization
(lowercase everything, sort everything, coerce everything to strings) fights
instability and manufactures puns; aggressive type preservation fights puns
and, done casually, reintroduces instability (embedding a language-specific
repr string preserves type at the cost of platform dependence). The rule set
below is best read as a sequence of decisions about exactly where to
normalize and exactly where to distinguish, with each decision justified
against both threats at once.

\section{The rule set}\label{ch05:rules}

\begin{ruleenv}[Canonical encoding, v0.1 profile]\label{ch05:rule:master}
The canonical evidence encoding:
\begin{enumerate}
\item is deterministic UTF-8 JSON with sorted string keys and minimal
separators, terminated by a single line feed;
\item rejects unordered containers (sets and frozen sets), binary floating-
point values, NaN, and infinities;
\item tags exact decimals, tuples, and enumeration members as typed objects,
so that $\mathrm{Decimal}(\texttt{"1.0"})$, the string \texttt{"1.0"}, and
the integer $1$ have pairwise distinct encodings, and reserves the tag key
so payload data cannot forge a tagged value;
\item gives each public type exactly one declared serialization route; and
\item is bounded and cycle-aware: it enforces declared depth, node-count,
byte-size, and integer-width limits, and fails closed on any reference
cycle.
\end{enumerate}
\end{ruleenv}

Before the clause-by-clause development, the profile's central property
can be stated once, in the vocabulary of
Definition~\ref{ch05:def:encoding}.

\begin{proposition}[Determinism and injectivity on the admitted domain]
\label{ch05:prop:injective}
On the domain of values the profile admits, the canonical encoding is a
function of value content alone---independent of construction history,
platform, process, and run---and is injective up to the profile's
declared identifications: two admitted values encode to the same bytes
if and only if they present the same canonical content under their
declared serialization routes.
\end{proposition}

\begin{designargument}
Determinism: every source of environmental variation is either
eliminated (sorted keys remove ordering dependence; minimal separators
remove formatting dependence; UTF-8 with direct non-ASCII emission
removes escaping choices) or excluded from the domain (unordered
containers, floats, and locale-sensitive text transformations are
rejected, so no admitted value's encoding consults them). Injectivity:
the JSON grammar is uniquely parseable; distinct admitted scalars
receive distinct token forms, with the JSON-inexpressible distinctions
(decimal, tuple, enumeration) carried by tagged objects under a key that
is reserved and therefore cannot arise from payload data; and route
uniqueness (Rule~\ref{ch05:rule:oneroute}) ensures no value has two
derivations. The argument is a design argument, checked by the
acceptance, rejection, and determinism fixture families of the bound
suite (Section~\ref{ch05:testing}); it is not a machine-checked proof,
and the open obligation to formalize it is recorded in Chapter 12.
\end{designargument}

Each clause now gets its own section, with the rationale, the worked
encodings, and the specific incident or attack it answers.

\section{Deterministic UTF-8 JSON with sorted keys}\label{ch05:json}

\subsection{The base grammar}

The encoding's surface grammar is JSON: objects with string keys, arrays,
strings, integers, booleans, and null. JSON is chosen not for expressiveness
---the whole point of the profile is to restrict it---but for auditability:
the encoded bytes are readable by a human reviewer and parseable by
independent tooling in any language, so a digest dispute can be settled by
inspection. Three determinism decisions turn the grammar into an encoding.

\emph{Key order.} Object keys are emitted in sorted order, so that the bytes
do not depend on insertion order, hash seeding, or any other artifact of how
the in-memory mapping was built. Two mappings with the same key--value pairs
encode identically regardless of construction history. The sort is over the
key strings themselves as sequences of code points; Section~\ref{ch05:jcs}
returns to the subtlety that this ordering choice is not the same as RFC
8785's.

\emph{Separators and whitespace.} The encoding uses minimal separators---a
comma between items, a colon between key and value, no spaces, no
indentation---so that formatting conventions cannot create spurious byte
differences. Pretty-printing is a display concern; display transformations
of evidence are computed from the canonical bytes, never the reverse.

\emph{Text encoding and terminator.} The bytes are UTF-8 with non-ASCII
characters emitted directly rather than escaped, and the encoding is
terminated by exactly one line feed. The terminator is a profile decision
worth flagging: it makes an encoded evidence record a well-formed
single-line text file on POSIX conventions, at the cost of a deliberate,
documented divergence from encoders that emit no terminator
(Section~\ref{ch05:jcs}).

\begin{example}[Worked base encoding]\label{ch05:ex:base}
The mapping with pairs \texttt{b} $\mapsto 2$, \texttt{a} $\mapsto 1$,
\texttt{c} $\mapsto$ [\texttt{true}, \texttt{null}] encodes as the bytes
\begin{center}
\texttt{\{"a":1,"b":2,"c":[true,null]\}}
\end{center}
followed by one line feed---regardless of the order in which the pairs were
inserted, the platform, or the process's hash seed. The same mapping built
by parsing, by literal construction, or by merging two mappings encodes to
the same twenty-nine bytes plus terminator, and therefore to the same
SHA-256 digest.
\end{example}

\subsection{Text admissibility}

Determinism of the byte layer is necessary but not sufficient: the
\emph{content} of strings can smuggle instability and confusion. The profile
therefore restricts admissible text before it restricts anything else. Every
string entering the encoding---keys and values alike---is validated to
exclude three Unicode categories: control characters, invisible format
controls, and surrogate code points. Control characters make encoded
evidence unsafe to display and diff; invisible format controls (zero-width
spaces, directional overrides) make two visually identical strings
byte-distinct, which is Threat~\ref{ch05:threat:instability} weaponized
against human review; and surrogates are not characters at all---their
presence means the value was manufactured through an encoding error, and
encoding them produces byte sequences that independent UTF-8 parsers
legitimately reject or repair differently.

Rejection, not repair, is the deliberate choice. A profile that silently
strips a directional override produces clean bytes from dirty input---and
thereby makes the dirty input and its cleaned twin encode identically, a pun
between a string that renders deceptively and one that does not. The
kernel's position is that deceptive text is not admissible evidence, and the
place to discover that is at encoding time, with an error naming the field.

\subsection{The hash layer}\label{ch05:hashlayer}

The encoding exists to be hashed, and the hash layer has its own small
profile, stated here once so the rest of the monograph can say
``canonical digest'' without ambiguity. The digest function is SHA-256
over the canonical encoding's bytes---terminator included---and nothing
else: no salt, no key, no truncation, no double hashing. The choice is
conservative and deliberately uninteresting: a single, universally
implemented, unbroken cryptographic hash, so that any party with any
toolchain can recompute a digest from bytes and so that a digest value
in an evidence record has exactly one possible derivation.

Digest \emph{values}, when they appear as fields in payloads and
records, are governed by the identifier discipline of
Section~\ref{ch05:identifiers} in its strictest form: exactly sixty-four
hexadecimal digits, carried in uppercase---the profile's declared case,
matching the reference layer's digest normalization
(Chapter~\ref{ch03:identity})---with case normalization applied at the
field boundary so that a digest arriving in lowercase is normalized on
intake rather than stored in mixed conventions. The alternative---tolerant, case-insensitive
digest comparison scattered across consumers---is a standing invitation
for one consumer to forget the tolerance, at which point the same
digest exists as two unequal strings on two sides of a comparison.

One caveat belongs to this layer and is stated here rather than
repeated: every identity claim of the form ``equal digest, therefore
equal bytes'' holds up to the collision resistance of SHA-256. The
kernel makes no contribution to, and no claims about, that
cryptographic property; it inherits the standard assumption and records
it as an assumption. What the kernel does ensure is the converse
direction and the preimage discipline around it: equal admitted values
always produce equal digests (Proposition~\ref{ch05:prop:injective}),
and every digest in a kernel record names bytes that are retained or
pinned somewhere a verifier can reach, so the assumption is only ever
invoked about bytes that exist, never about bytes taken on faith.

\section{The rejection set: what cannot be evidence}\label{ch05:rejections}

Clause (2) of Rule~\ref{ch05:rule:master} names four rejections. Each is a
value family that mainstream serializers accept, and each is rejected here
because it is \emph{poisonous to digests}: its presence anywhere in a payload
makes the encoding either nondeterministic or many-to-one. This section
takes them in turn.

\subsection{Unordered containers}\label{ch05:sets}

A set has no canonical element order, so any encoding of a set must choose
one. Every available choice is bad. Iteration order is the worst: in the
host language it depends on insertion history and, for strings, on a
per-process hash seed, so the same set encodes differently in the same
program across runs---pure Threat~\ref{ch05:threat:instability}. Sorting the
elements repairs determinism for homogeneous sets of totally ordered
elements and fails for everything else: mixed-type sets have no total order,
and imposing an artificial one (sort by type name, then value) builds a
subtle language-specific convention into the byte layer, which is
instability displaced, not removed. Worse, sort-based set encoding
manufactures a pun: the set $\{1, 2\}$ and the \emph{list} $[1, 2]$ become
byte-identical, though one is an unordered membership assertion and the
other an ordered sequence---evidence with different semantics.

The kernel refuses the choice entirely: sets and frozen sets are rejected
with a typed error. A caller holding set-shaped data must decide, at the
call site and visibly in code review, what ordered artifact the evidence
actually is---usually a sorted list under an ordering the caller can state
and defend. The decision is forced to the layer that has the semantics,
instead of being made silently and identically-for-everyone by the encoder,
which is the one layer guaranteed not to know what the set means.

\subsection{Binary floats}\label{ch05:floats}

Binary floating-point values are rejected outright---not merely NaN and
infinities, but every float, including innocent-looking ones. The rationale
has three independent layers, any one of which would suffice.

First, representation: most decimal literals are not representable in
binary floating point. The value written \texttt{0.1} is stored as the
nearest double, $0.1000000000000000055511151231257827\ldots$; the value
written \texttt{0.3} is not the sum of the stored \texttt{0.1} and
\texttt{0.2}. Evidence records exist to carry exact commitments, and a float
is an inexact commitment wearing exact-looking digits.

Second, formatting: turning a double back into digits is a rendering
algorithm with a history of divergence---shortest-round-trip rendering
versus fixed precision, differences across languages, libraries, and
versions. Two systems holding bit-identical doubles can emit different
digit strings, and two different doubles can emit the same short string
under loose formatting. That is both threats of
Section~\ref{ch05:threats} at once, seated in the innermost byte layer.

Third, arithmetic provenance: floats invite computation before encoding,
and float computation is order-sensitive ($ (a+b)+c \neq a+(b+c) $ in
general), platform-tunable, and non-reproducible by default. A digest over
a float-bearing payload is a digest over the accumulated rounding history
of whatever machine produced it.

The kernel's alternative is Section~\ref{ch05:decimal}'s tagged exact
decimals: values that carry precisely the digits the recorder committed to,
no more and no fewer. The rejection of floats is enforced by exact type
check, which has a small but important corollary in the host language:
booleans, which the language implements as an integer subtype, are tested
for \emph{before} integers, so \texttt{true} encodes as a boolean and can
never be punned with the integer $1$---a miniature instance of the same
discipline.

\subsection{NaN}\label{ch05:nan}

NaN would be rejected even if floats were not, and it is worth recording
why, because NaN is uniquely corrosive to an evidence system. NaN is not a
value; it is the propagating absence of one. It is the only widely deployed
datum that is unequal to itself, so its presence breaks the reflexivity that
every comparison-based mechanism in this monograph assumes: a record
containing NaN cannot even be said to equal \emph{itself}, and ``same
evidence, same digest'' loses meaning at the semantic layer no matter what
the byte layer does. Serialization compounds this: the JSON grammar has no
NaN token, so encoders that accept NaN emit a nonstandard extension token
that independent parsers variously accept, reject, or translate---meaning
the encoded bytes are not, in fact, JSON, and cross-tool audit fails on
exactly the payloads containing the least meaningful data. A NaN in an
evidence payload is always a symptom: some upstream computation failed and
the failure was allowed to flow. The kernel's fail-closed contract requires
surfacing that failure at its origin, not laundering it into a hashed
record whose digest then faithfully commits to nothing.

\subsection{Infinities}\label{ch05:inf}

Positive and negative infinity share NaN's grammatical problem---no JSON
token, nonstandard extensions, parser divergence---without its reflexivity
pathology, and they add a subtler one: they encode a claim (``this quantity
exceeds every finite bound'') that is essentially never the fact a recorder
means to commit to. In practice an infinity in a payload records an
overflow, a division by zero, or a sentinel convention (``no limit'')
that belongs in the schema as an explicit optional field, not in the
number line. Because the kernel's exact-decimal carrier is finite by
construction---the decimal type admits non-finite values, and the encoder
checks finiteness and rejects them---the profile makes ``every encoded
number is a finite, exact decimal or a bounded integer'' an invariant a
consumer can rely on without reading any schema.

\begin{invariant}[Finite exact numerics]\label{ch05:inv:finite}
Every numeric value in a canonical evidence payload is either an integer
within the declared width bound or a finite tagged decimal. No encoded
payload contains a binary float, a NaN, or an infinity; fixtures for each
rejection are part of the bound suite.
\end{invariant}

\section{Tagged exact values}\label{ch05:decimal}

\subsection{The pun problem, concretely}

Clause (3) is the profile's answer to Threat~\ref{ch05:threat:puns}, and the
motivating instance is numeric. Consider three values a program might hold:
the exact decimal constructed from the digits \texttt{1.0}; the five-
character string \texttt{"1.0"}; and the integer $1$. A naive encoder maps
the decimal to the JSON number token \texttt{1.0}, the string to
\texttt{"1.0"}, and the integer to \texttt{1}---three distinct byte
strings, which looks fine until one notices that the decimal's token
collides with what a float encoder would emit, that many decimal encoders
``simplify'' \texttt{1.0} to \texttt{1} and collide with the integer, and
that string-typed numeric fields (ubiquitous in ledger data) collide with
stringified decimals. Across real serializer configurations, some pair of
these three values collides somewhere. An evidence system cannot tolerate
that: a quantity asserted as an exact decimal with one decimal place of
recorded precision is a different commitment from the bare integer and
from a piece of prose that happens to contain digits.

\subsection{The tagging scheme}

The kernel encodes every type-distinct scalar that JSON cannot natively
distinguish as a small tagged object under a reserved key. The tag key is
\texttt{\$cas\$type}; its value names the carried type; the remaining fields
carry the value in a deterministic rendering.

\begin{example}[The three encodings differ]\label{ch05:ex:decimal}
Under the v0.1 profile:
\begin{itemize}
\item the exact decimal $\mathrm{Decimal}(\texttt{"1.0"})$ encodes as
\texttt{\{"\$cas\$type":"decimal","value":"1.0"\}};
\item the string \texttt{"1.0"} encodes as \texttt{"1.0"};
\item the integer $1$ encodes as \texttt{1}.
\end{itemize}
The three byte strings are pairwise distinct, so the three digests are
pairwise distinct (up to hash collision), and no evidence about one value
can validate a record about another. The decimal's digits are rendered in
fixed-point form from the decimal's own exact state---never through a
binary float---so the rendering is a function of the committed digits
alone.
\end{example}

\begin{remark}[Precision is part of the value]\label{ch05:rem:precision}
Because the rendering preserves the decimal's stored exponent,
$\mathrm{Decimal}(\texttt{"1.0"})$ and $\mathrm{Decimal}(\texttt{"1.00"})$
encode differently: \texttt{"1.0"} versus \texttt{"1.00"} in the tagged
value field. This is deliberate. In measurement contexts the recorded
precision is part of the assertion---a companion candidate records
calibrated quantities where one decimal place versus two is a different
claim about the instrument. Callers for whom trailing zeros are \emph{not}
significant must normalize before encoding, visibly, at the call site; the
encoder will not decide for them. The same one-way principle as
Section~\ref{ch05:sets}: normalization is a semantic act and belongs to
the layer that knows the semantics.
\end{remark}

The same scheme carries two further distinctions. Tuples---used in the
kernel for fixed-shape positional data---encode as
\texttt{\{"\$cas\$type":"tuple","items":[\ldots]\}}, distinct from JSON
arrays, which carry variable-length sequence data; a loader can therefore
reject a sequence where a fixed shape was committed, and vice versa.
Enumeration members encode as a tagged object carrying both the
enumeration's fully qualified type name and the underlying value, so the
member \textsc{open} of the admission-status enumeration can never be
punned with the bare string \texttt{"OPEN"}---which matters directly for
Chapter 4's boundary carrier, whose digested form must not be forgeable
from string-typed data that happens to spell the same words.

\begin{example}[Worked enumeration encoding]\label{ch05:ex:enum}
Let \textsc{open} be a member of an admission-status enumeration whose
fully qualified type name is
\texttt{cas\_core\_v8.semantics.AdmissionStatus} and whose underlying
value is the string \texttt{"OPEN"}. Encoded directly, the member
produces (with keys shown sorted, as emitted):
\begin{center}
\texttt{\{"\$cas\$type":"enum",}\\
\texttt{"enum\_type":"cas\_core\_v8.semantics.AdmissionStatus",}\\
\texttt{"value":"OPEN"\}}
\end{center}
while the bare string \texttt{"OPEN"} produces the six bytes
\texttt{"OPEN"} (with quotes). The tagged form binds three facts the
bare string cannot: that the value is an enumeration member at all,
\emph{which} enumeration it belongs to, and what its underlying value
is. Two enumerations that happen to share the word \textsc{open}---an
admission status and, say, a disclosure status---therefore encode
distinctly, and a digest over a payload containing one can never
validate a payload containing the other. The enumeration check runs
before the scalar checks precisely because scalar-backed enumerations
would otherwise decay to their underlying values and lose all three
facts. Note also the dual reading: a \emph{declared view} may choose to
project an enumeration to its plain string (Chapter 4's boundary view
does exactly this), and under Rule~\ref{ch05:rule:oneroute} that
projection is stable and unambiguous because it is the type's only
route; what the tagging forbids is not projection but \emph{collision}
---the projected string lives inside a view whose shape is fixed, not
loose in a position where a raw enumeration member could also appear.
\end{example}

\subsection{Reserving the tag}

A tagging scheme creates its own forgery surface: if payload data may
contain the tag key, then a mapping in user data with key
\texttt{\$cas\$type} and value \texttt{"decimal"} is byte-
indistinguishable from an encoder-produced tagged decimal, and the pun
returns one level up. The profile closes this by reservation: the tag key
is forbidden as an ordinary mapping key anywhere in admissible payloads,
enforced at encoding time (and independently at registry-load time,
Chapter 6). Within the admitted domain, an encoded object bearing the tag
key was produced by the encoder's own tagging path---there is no other
way for those bytes to arise---which is what makes clause (ii) of
Definition~\ref{ch05:def:encoding} (injectivity) hold across the tagged
extension: the encoding remains one-to-one because the tagged and untagged
worlds cannot collide by construction.

\section{One serialization route per public type}\label{ch05:oneroute}

Clause (4) is a rule about the codebase rather than about any single byte
string, and it exists because multi-route encoding is how determinism dies
in mature systems even when every individual route is deterministic.

Consider a public type---say, the artifact-identity record of
Section~\ref{ch05:byteid}---that can reach bytes two ways: a hand-written
view that emits selected fields in a chosen shape, and a generic reflective
path that walks the record's fields mechanically. Both are deterministic.
They are also different: the hand-written view may include a derived
field the reflective path omits, order or name fields differently, or
render a nested value through a different sub-route. Every digest computed
on route A is incomparable with every digest computed on route B, and
\emph{nothing fails}: each route round-trips its own output, each verifies
its own digests, and the system contains two parallel, internally
consistent evidence universes. The divergence surfaces only when a record
digested on one route is verified on the other---typically at the worst
seam, between two components maintained by different people, each of whom
is provably ``correct.''

\begin{ruleenv}[Single-route encoding]\label{ch05:rule:oneroute}
Each public type has exactly one declared serialization route into the
canonical encoding. Types with a declared evidence view are encoded through
that view and only that view; the generic structural path is a fallback for
types without one, and the encoder consults the declared view first, so no
admissible value has two derivations. Adding a second route to an existing
type is a breaking change to evidence comparability and is treated as such.
\end{ruleenv}

The v0.1 encoder implements the rule by a fixed dispatch order: tagged
scalars first, then declared views, then the generic structural walk for
plain field-record types, with the reflective walk deliberately reimplemented
inside the bounded traversal of Section~\ref{ch05:bounded} rather than
delegated to a library helper whose recursion is neither cycle- nor
depth-bounded. The rule's payoff is a property worth stating:

\begin{proposition}[Route uniqueness implies comparability]
\label{ch05:prop:comparability}
Under Rule~\ref{ch05:rule:oneroute}, for any two encoded payloads of the
same public type, digest inequality implies value distinction under the
type's declared view, and digest equality implies byte-identical canonical
content. Without route uniqueness, neither implication holds: digest
inequality may reflect route choice, and digest equality on one route says
nothing about encodings on another.
\end{proposition}

\subsection{A worked end-to-end encoding}\label{ch05:worked}

The rules compose; a full example shows how. Consider the Ring-4
research boundary of Chapter 4, whose declared view presents four
fields: the admission ring (integer), the admission status and license
status (enumerations, projected by the view to their plain string
values), and the disclosure flag (boolean). The encoder consults the
declared view (Rule~\ref{ch05:rule:oneroute}), obtains a four-key
mapping, validates each key and each string value as admissible text,
sorts the keys, and emits:
\begin{center}
\texttt{\{"admission\_ring":4,"license\_status":"NOT LICENSED",}\\
\texttt{"public\_disclosure\_authorized":false,"status":"OPEN"\}}
\end{center}
followed by the terminating line feed. Every profile decision is
visible in these bytes. The keys appear in sorted order regardless of
the view's construction order. The ring is a bare JSON integer---
admissible because bounded---and could not have been a float; had the
view held the ring as a float $4.0$, encoding would have failed with a
typed error rather than emitting \texttt{4} or \texttt{4.0}. The two
status strings are the view's declared projections; they are ordinary
strings here because the view says so, and the same enumeration members
encoded \emph{outside} this view would take the tagged form of
Example~\ref{ch05:ex:enum}---no ambiguity arises because each public
type has one route and the route fixes the shape. The boolean is
\texttt{false}, not $0$: exact-type dispatch keeps the boolean/integer
distinction. There are no spaces after separators, no escaping of the
space inside \texttt{"NOT LICENSED"} (UTF-8 direct emission), and
exactly one terminator byte at the end. The SHA-256 of these bytes is
the carrier's canonical digest; by
Proposition~\ref{ch05:prop:injective} any process on any platform that
holds an equal boundary value computes the same digest, and no
admissible value other than this boundary's view content computes it
(up to hash collision, per Section~\ref{ch05:hashlayer}).

The example also illustrates the division of labor between this chapter
and Chapter 4: the encoding guarantees that these bytes faithfully and
uniquely name the carrier's content; it is the boundary discipline, not
the encoding, that determines what the content \emph{means}---and, per
the standing invariants, the digest of these bytes evidences their
integrity, never anyone's authorization of them.

\section{Bounded, cycle-aware canonicalization}\label{ch05:bounded}

Clause (5) addresses a family of failures that are not about which bytes
come out but about whether the encoder terminates and stays within resource
bounds when handed adversarial or accidentally pathological structure. An
encoder that is part of a fail-closed evidence path is itself an attack
surface: if hashing a payload can be made to hang, recurse without bound, or
allocate without bound, then the evidence path is a denial-of-service lever
against every operation that must hash before proceeding---which, in this
kernel, is every operation.

The v0.1 profile enforces four declared limits and one structural check.

\emph{Depth.} Traversal beyond a declared nesting depth (64 in v0.1) fails
closed. Legitimate evidence in this program is shallow; a payload nested
hundreds deep is either a bug or a crafted recursion bomb, and the encoder
refuses both identically.

\emph{Node count.} Traversal visits at most a declared number of nodes
(10{,}000 in v0.1). This bounds work even for wide, shallow structures that
evade the depth limit.

\emph{Byte size.} Accumulated text is metered during traversal and the final
encoding is checked against a declared output bound (8 MiB in v0.1), so a
payload cannot balloon past what downstream storage and review can handle.
The in-flight metering matters: it stops runaway text at the point of
accumulation instead of after a gigabyte-scale join.

\emph{Integer width.} Integers beyond a declared bit width (2{,}048 bits in
v0.1) are rejected. Arbitrary-precision integers are otherwise a quiet
amplification channel---a small in-memory object whose decimal rendering is
enormous---and no legitimate evidence quantity in this program approaches
the bound.

\emph{Cycles.} The traversal maintains the set of composite objects on the
current path; revisiting an active object fails closed immediately with a
cycle error. The alternative behaviors are both unacceptable: infinite
recursion is a crash at best, and silent cycle-breaking (encode the first
visit, elide the second) makes the encoding depend on traversal entry
point---two references to the same cyclic structure would encode
differently depending on where the walk began, which is
Threat~\ref{ch05:threat:instability} in structural clothing. Note the check
is scoped to the \emph{active path}: acyclic sharing (the same immutable
value referenced twice in different branches) is legal and encodes by
value at each occurrence, which is the correct evidence semantics---the
payload commits to content, not to in-memory object graph shape.

The four values themselves are engineering margins, not measured
quantities: the governing sources record no payload census from which
they were derived. The honest statement is that each bound sits far
above every payload the bound suites exercise while remaining cheap to
enforce, and that a legitimate payload ever approaching a bound is, by
Remark~\ref{ch05:rem:limits}, a forced profile-version event rather
than a quiet tuning change.

\begin{remark}[Limits are part of the profile]\label{ch05:rem:limits}
The four numeric limits are declared constants of the profile, not tuning
parameters: two implementations of the encoding agree only if they enforce
the same bounds, since a payload accepted by one and rejected by the other
breaks fail-closed interoperability exactly at the margin. Any future
revision of the limits is therefore a version change of the encoding
profile itself, visible in the profile identity that receipts record.
\end{remark}

\section{Byte identity versus comparison identity}\label{ch05:byteid}

\subsection{Two identities, both recorded}

Everything so far concerns values the kernel encodes. A second identity
problem attaches to artifacts the kernel merely \emph{reads}: source files,
ledgers, schemas, compiler bytes. For these, the bytes already exist, and
the question is what ``the same file'' means across platforms whose text
conventions differ. The program's operating history makes this concrete:
files authored on one platform arrive on another with carriage-return/
line-feed sequences rewritten to bare line feeds or vice versa---by
version-control autocrlf settings, editor defaults, or transfer tooling.
After such a rewrite the file's raw digest has changed while, on the
program's declared comparison convention, its content has not.

Collapsing this situation in either direction loses information that
incidents have shown to matter. Recording only the raw digest makes every
platform transit look like tampering, flooding review with false
mismatches until reviewers stop looking---the desensitization failure of
Threat~\ref{ch05:threat:instability}. Recording only a normalized digest
makes genuine byte differences invisible whenever they hide in line
endings, and byte-exactness is precisely what Chapter 6's registries and
Chapter 7's compiler verification pin.

\begin{ruleenv}[Dual identity]\label{ch05:rule:dual}
Every artifact identity records both a \emph{raw} identity---digest and
byte count of the exact bytes---and a declared \emph{comparison}
identity---digest and byte count after the single declared normalization,
CRLF-to-LF, and no other transformation. The two are separate fields with
separate roles; neither substitutes for the other, and no kernel operation
silently accepts one where the other is required.
\end{ruleenv}

The normalization is deliberately minimal and closed: replace each
two-byte CRLF sequence with LF, touch nothing else. It is not trailing-
whitespace trimming, not encoding transcoding, not Unicode normalization.
A one-transformation comparison identity can be reasoned about exactly;
an open-ended ``semantic'' normalization cannot.

\subsection{The line-ending census}

Alongside the two digests, the identity records a three-way census of the
artifact's line-ending sequences: the count of CRLF sequences, of lone
line feeds, and of lone carriage returns. The three counts are never
conflated into a single ``line count,'' because their joint pattern is
diagnostic in a way no aggregate is. All-CRLF with zero lone characters is
a consistent foreign-convention file; all-LF is a consistent native one;
nonzero counts in more than one column mean a \emph{mixed} file---almost
always the fingerprint of a partial rewrite, a bad merge, or an editor
that converted only the lines it touched; and lone carriage returns,
rare on modern platforms, flag either ancient tooling or bytes that were
never text. The identity derives a summary profile (CRLF, LF, CR, mixed,
or no line endings) from the census for human display, but the counts
themselves are the recorded evidence.

The census also underwrites an internal consistency check that turns the
dual identity from two assertions into one verifiable structure: since
the only transformation is CRLF-to-LF and each application shortens the
artifact by exactly one byte, the normalized byte count must equal the
raw byte count minus the CRLF count. An identity record violating this
arithmetic is rejected at construction---it could not have been produced
by measuring any real byte string, so it is either corrupt or fabricated.

\begin{proposition}[Census arithmetic]\label{ch05:prop:census}
For any byte string, the CRLF-to-LF normalized length equals the raw
length minus the number of CRLF sequences, and the total of
$2\cdot(\textrm{CRLF count}) + (\textrm{lone LF count}) + (\textrm{lone
CR count})$ cannot exceed the raw length. Both facts are enforced as
construction-time checks on identity records, so a record's five numeric
fields are mutually corroborating rather than independently assertable.
\end{proposition}

\subsection{The incidents this answers}

Three recurring incident shapes, all previously observed in the program's
operations, are each made visible-instead-of-silent by the dual identity.
A cross-platform checkout changes raw digests en masse while comparison
digests hold: the pattern ``raw differs, LF-normalized matches, CRLF
census changed'' names the event exactly, and review can wave it through
\emph{as that event} rather than as an unexplained mismatch. A genuine
content change hidden in a reformatting commit shows the opposite
signature---comparison digests differ---which no amount of newline noise
can mask. And a partial conversion produces a mixed census where either
consistent convention stood before, surfacing the half-rewritten file
that per-digest comparison alone would report only as ``different,''
without the why.

\section{Identifier profiles}\label{ch05:identifiers}

Text in evidence records serves two masters. Machine identifiers---record
identities, namespace names, version strings, hash fields---must be
comparable, greppable, unambiguous in logs, and safe across every tool in
the audit chain. Prose---titles, notes, definitions quoted from source
documents---must be carried faithfully, including in languages and scripts
the identifiers' rules would forbid. A single admissibility rule for both
either cripples prose or opens identifiers to confusable-character and
normalization attacks. The profile therefore declares two tiers.

\begin{ruleenv}[Two-tier text]\label{ch05:rule:tiers}
Review-visible machine identifiers satisfy a declared restricted profile:
nonblank, trimmed, ASCII-only, drawn from an explicit identifier alphabet
(alphanumerics plus a short list of structural punctuation), and a fixed
point of compatibility normalization---the identifier must already equal
its own NFKC normal form. Prose fields carry unrestricted Unicode subject
only to the global text admissibility of Section~\ref{ch05:json}
(no controls, no invisible format characters, no surrogates). Prose never
participates in identifier comparison, and identifiers are compared as
exact byte sequences with no case folding or locale participation.
\end{ruleenv}

Each clause of the restricted profile closes a specific hole. ASCII-only
removes the entire homoglyph space---no Cyrillic \texttt{a} posing as
Latin \texttt{a} in a record identity. The NFKC fixed-point requirement is
subtler: rather than normalizing identifiers on intake (which would make
two visually distinct inputs collide post-hoc, a pun manufactured by the
kernel itself), the profile requires that the identifier \emph{already be}
normal, rejecting anything that normalization would change. The kernel
thus never rewrites an identifier; it only refuses non-canonical ones---
the same reject-don't-repair stance as Section~\ref{ch05:json}, applied
one tier up. Hash-valued fields get a stricter micro-profile still: exactly
sixty-four hexadecimal digits, normalized to uppercase, the declared
case, at the boundary, so that digest comparison is byte comparison
everywhere downstream and a mixed-case digest can never split into two
apparently different values.

\section{Comparison with RFC 8785 (JCS)}\label{ch05:jcs}

The obvious question about any bespoke canonical JSON profile is why it is
not simply the standard one. RFC 8785, the JSON Canonicalization Scheme
\cite{jcs}, solves an overlapping problem: deterministic bytes for JSON
data, for exactly the hashing-and-signing use cases this chapter cares
about. The kernel's profile agrees with JCS on the load-bearing
architecture and departs from it in documented, deliberate ways; this
section states both sides precisely, because ``roughly JCS'' is exactly
the kind of claim the monograph's discipline forbids.

\subsection{Where the profile follows JCS}

The agreements are the structural ones. Both use UTF-8 as the sole byte
encoding. Both emit objects with sorted keys, so member order carries no
information. Both use minimal separators with no insignificant whitespace.
Both emit non-ASCII text directly rather than ASCII-escaped. Both are
designed so that hashing the output is meaningful: same data, same bytes,
same digest. A reader who knows JCS can read the kernel's encoded evidence
without relearning anything at the surface level.

\subsection{Where the profile is stricter}

Every departure is a restriction or an extension motivated by the evidence
threat model, and each is flagged in Table~\ref{ch05:tab:jcs}.

\emph{Numbers.} The deepest divergence. JCS inherits I-JSON's number
model: numbers are IEEE-754 doubles, serialized by the ECMAScript
shortest-round-trip algorithm. This is the right call for JCS's goal---
canonicalizing arbitrary interoperable JSON---and it is unacceptable
here, for every reason in Section~\ref{ch05:floats}: doubles are inexact
carriers, and shortest-round-trip rendering, while deterministic, commits
evidence to binary-rounding artifacts. The kernel rejects floats entirely
and carries exact quantities as tagged decimals. In JCS terms: the
kernel's numeric domain is integers (bounded) plus a tagged-object
extension, and no value in the kernel's domain exercises JCS's
double-serialization machinery at all.

\emph{Key ordering.} JCS sorts object members by the UTF-16 code units of
their names; the kernel's encoder sorts by Unicode code points. The two
orders agree on all-ASCII keys---which covers every machine identifier
under Rule~\ref{ch05:rule:tiers}---but can differ when keys mix
supplementary-plane characters with code points in the upper basic
multilingual plane, because surrogate-pair code units interleave
differently than code points. The kernel documents this as a divergence
rather than hiding it: its encoded bytes for such keys are \emph{not}
JCS bytes. The practical exposure is narrow (non-ASCII keys are rare in
kernel payloads and forbidden in identifier positions), but narrow is not
zero, and the profile says so. The governing sources do not record why
code-point order was preferred over adopting JCS's code-unit sort
outright; the natural reading, offered as reconstruction rather than
record, is that code-point comparison is the host language's native
string order, so the canonical sort needs no hand-written UTF-16
transcoding step---one fewer transformation in exactly the place where
this chapter counts transformations as attack surface---and
byte-for-byte JCS compatibility was already forfeited by the terminator
and the tagged extension.

\emph{Type discrimination.} JCS canonicalizes JSON's data model and
therefore cannot distinguish what JSON cannot: it has no tuples, no
decimals-as-distinct-from-numbers, no enumerations. The kernel's tagged
extension (Section~\ref{ch05:decimal}) is strictly outside JCS's scope---
JCS would happily canonicalize the tagged objects as ordinary objects,
which is compatible but meaningless without the reservation rule that
makes tags unforgeable.

\emph{Admissibility.} JCS canonicalizes whatever valid JSON it is given;
it has no rejection set. The kernel rejects sets and floats at the host-
language boundary, rejects control and format characters inside strings,
rejects the reserved tag key in data, and enforces the resource bounds of
Section~\ref{ch05:bounded}, none of which JCS contemplates. In the same
spirit, the kernel's registry layer rejects duplicate keys at parse time
(Chapter 6), where JCS---operating downstream of a parser---assumes
I-JSON well-formedness rather than enforcing it.

\emph{Output framing.} The kernel terminates encodings with a line feed;
JCS output is unterminated. A JCS digest and a kernel digest of the same
object therefore differ even when every byte before the terminator
agrees. This is the one divergence with no semantic content, retained
because the kernel's encoded records live as files in a line-oriented
audit toolchain; it is called out so that no consumer expects cross-
verification against a JCS implementation to succeed byte-for-byte.

\begin{table}[ht]
\centering
\begin{tabular}{lll}
\toprule
Dimension & RFC 8785 (JCS) & CAS Core 8 v0.1 profile \\
\midrule
Byte encoding & UTF-8 & UTF-8 (same) \\
Key order & UTF-16 code units & code points (differs off-ASCII) \\
Whitespace & none & none (same) \\
Numbers & IEEE doubles, ES6 rendering & floats rejected; tagged decimals \\
NaN / infinity & inexpressible in I-JSON & rejected with typed error \\
Sets / tuples / enums & out of scope & rejected / tagged / tagged \\
String admissibility & any JSON string & controls, format chars, surrogates rejected \\
Duplicate keys & assumed absent (I-JSON) & rejected at parse (registry layer) \\
Resource bounds & none & depth, nodes, bytes, integer width \\
Terminator & none & single LF \\
\bottomrule
\end{tabular}
\caption{Agreement and divergence between the JCS canonicalization and the
kernel's evidence-encoding profile. ``Rejected'' means a typed error at
encoding time; the kernel is a restricted profile with a tagged extension,
not a JCS implementation.}
\label{ch05:tab:jcs}
\end{table}

\begin{remark}[Relation to I-JSON]\label{ch05:rem:ijson}
JCS is defined over the I-JSON message profile---the interoperable
subset of JSON that excludes duplicate keys, surrogates, and numbers
outside double precision \cite{jcs}. The kernel's admitted domain is,
by construction, inside I-JSON's spirit and stricter than its letter:
duplicate keys are actively rejected rather than assumed absent,
surrogate and format-control code points are rejected rather than
discouraged, and the number question is mooted by the float rejection.
One divergence runs the other way and should be flagged: I-JSON bounds
integers to the double-representable range, while the kernel admits
integers up to its declared 2{,}048-bit width, carried as exact JSON
integer tokens. A consumer parsing kernel evidence with a
double-backed JSON library would silently round such integers; the
kernel's own parse layer does not, and the profile documents that
evidence consumers must use exact-integer parsing. In practice kernel
payloads carry digests as strings and quantities as tagged decimals,
so wide integers are rare---but rare is not absent, and the profile
says which side of the divergence it stands on.
\end{remark}

The summary judgment: the kernel follows JCS wherever JCS's choice is
about determinism, and overrides JCS wherever JCS's choice is about
\emph{tolerance}---of floats, of arbitrary strings, of unbounded
structure---because an evidence kernel's job is precisely to be
intolerant there. The profile should be read as ``JCS's architecture,
minus JCS's number model, plus a typed extension and a rejection set,''
and receipts that name the encoding profile name the kernel's, never
JCS, so no consumer can mistake one for the other.

\section{A failure-mode gallery}\label{ch05:gallery}

The chapter closes with a gallery of concrete failures, each stated as it
would occur in an undisciplined encoder, with the clause of the profile
that blocks it. The gallery serves the same role as Chapter 4's misuse
scenarios: it is the rule set exercised against realistic trouble, and
each entry corresponds to a rejection or determinism fixture family in
the bound suite.

\begin{example}[Hash-seed shuffle]\label{ch05:ex:seed}
A payload containing a set of identifier strings encodes in iteration
order; the process's string-hash seed differs per run; the same evidence
hashes differently on every execution, and replay verification fails
intermittently in a way that vanishes under debugging (fixed seeds).
\emph{Blocked by} the set rejection of Section~\ref{ch05:sets}: the
payload never encodes; the caller must supply a sorted list.
\end{example}

\begin{example}[The point-one trap]\label{ch05:ex:pointone}
A measurement of \texttt{0.1} units is stored as a double and encoded as
a number. A second system recomputes the quantity as
\texttt{0.30000000000000004}-style arithmetic residue, or renders the
same double with different formatting, and the digests diverge---or,
under loose formatting, two genuinely different stored quantities render
to the same short string and the digests collide. \emph{Blocked by} the
float rejection of Section~\ref{ch05:floats} and the tagged-decimal
carrier: the recorder must commit to exact digits, and those digits are
what is hashed.
\end{example}

\begin{example}[NaN equals nothing]\label{ch05:ex:nan}
An upstream failure yields NaN in a numeric field; the encoder emits a
nonstandard token; one downstream parser errors, another accepts, a third
reads null. Meanwhile the record is unequal to itself under the very
comparison semantics the registry relies on. \emph{Blocked by} the NaN
rejection of Section~\ref{ch05:nan}: the upstream failure surfaces at
encoding time as a typed error naming the field.
\end{example}

\begin{example}[Decimal--string--integer pun]\label{ch05:ex:pun}
A ledger field holding the string \texttt{"1.0"} is compared, via digest,
against a record whose author stored the decimal $1.0$; an untagged
encoder emits identical bytes for both, and evidence about a prose field
validates a numeric commitment. \emph{Blocked by} the tagging of
Section~\ref{ch05:decimal}: the encodings differ
(Example~\ref{ch05:ex:decimal}).
\end{example}

\begin{example}[Tag forgery]\label{ch05:ex:tagforgery}
An author embeds \texttt{\{"\$cas\$type":"decimal","value":"2.0"\}} as
ordinary data in a payload, crafting bytes identical to an encoder-
produced tagged decimal. \emph{Blocked by} the reservation rule: the tag
key is rejected in data positions at encoding time and again at registry
load, so the forged bytes are not producible through the admitted path.
\end{example}

\begin{example}[Two-route drift]\label{ch05:ex:tworoutes}
A convenience helper serializes a public record type reflectively for a
cache, while the evidence path uses the type's declared view; a field is
added to the view only; cache digests and evidence digests silently fork,
and a verification months later compares digests from different
universes. \emph{Blocked by} Rule~\ref{ch05:rule:oneroute}: one route per
public type, with route addition treated as a breaking change.
\end{example}

\begin{example}[Autocrlf masquerade]\label{ch05:ex:autocrlf}
A checkout rewrites a pinned source file from LF to CRLF; raw digests
change across the whole tree; a reviewer, habituated to platform noise,
approves a diff in which one file's change is \emph{not} line-ending
noise. \emph{Blocked by} the dual identity and census of
Section~\ref{ch05:byteid}: the noise files show ``raw differs, LF
matches, census flipped''; the real change shows ``LF differs,'' a
signature no newline rewrite can produce.
\end{example}

\begin{example}[Invisible ink]\label{ch05:ex:invisible}
An identifier contains a zero-width character; it renders identically to
its clean twin in every review surface, while comparing unequal and
hashing differently---two records that every human reads as one.
\emph{Blocked by} text admissibility (format controls rejected) and the
ASCII identifier profile of Section~\ref{ch05:identifiers}.
\end{example}

\begin{example}[Compatibility confusable]\label{ch05:ex:confusable}
An identifier uses a Unicode compatibility character (a full-width Latin
letter, a ligature) that NFKC would fold to the plain form; depending on
which layer normalizes, the identifier intermittently equals its folded
twin. \emph{Blocked by} the NFKC fixed-point requirement: the non-normal
form is rejected at intake, and the kernel itself never folds, so there
is no layer-dependent equality.
\end{example}

\begin{example}[Recursion bomb and cycle]\label{ch05:ex:bomb}
A crafted payload nests two thousand levels deep, or contains a mapping
that transitively contains itself; a naive encoder overflows the stack or
recurses forever inside the evidence path, denying service to every
operation that hashes first. \emph{Blocked by} the bounds and active-path
cycle check of Section~\ref{ch05:bounded}: depth and node limits fail
closed with typed errors; the cycle is detected at first re-entry.
\end{example}

\begin{example}[Big-integer balloon]\label{ch05:ex:bigint}
A payload carries $2^{100000}$ in a field where no legitimate quantity
exceeds a few hundred bits; the decimal rendering alone is tens of
kilobytes, and a stream of such records exhausts storage metered by
record count. \emph{Blocked by} the integer-width bound; quantities that
genuinely need arbitrary width must be redesigned as explicit structures,
visibly, rather than smuggled through the number path.
\end{example}

\begin{example}[Locale case surprise]\label{ch05:ex:locale}
A pipeline ``normalizes'' identifiers by lowercasing them, and one
stage runs under a locale in which case mapping differs from the ASCII
expectation (the classic instance maps the letter \texttt{I} outside
the ASCII pair), so the same identifier lowercases to different byte
strings on different hosts. \emph{Blocked by} the identifier profile of
Section~\ref{ch05:identifiers}: identifiers are compared as exact bytes
with no case folding anywhere in the kernel, and the one sanctioned
case normalization---hexadecimal digest case---is locale-independent by
construction. A transformation that is not in the profile does not
happen; a transformation that is in the profile is total and
deterministic.
\end{example}

\begin{example}[Escape-style fork]\label{ch05:ex:escape}
Two encoders serialize the same non-ASCII string, one emitting raw
UTF-8 and the other emitting ASCII escape sequences; both outputs parse
back to the same string, but the bytes---and therefore the digests---
differ, and a digest computed by one toolchain fails verification in
the other. \emph{Blocked by} profile fixity: the canonical encoding
emits non-ASCII directly, always, and receipts name the encoding
profile, so a consumer either implements the declared emission or is
not computing canonical digests at all. Parse-equivalence is explicitly
not the profile's identity relation; byte equality is.
\end{example}

\begin{example}[Terminator mismatch]\label{ch05:ex:terminator}
An independent implementation reproduces every rule but omits the final
line feed; all of its digests differ from the kernel's on every
payload, and the divergence is invisible to any inspection of the
pre-terminator content. \emph{Blocked by} documentation and fixtures
rather than by code: the terminator is a declared clause of
Rule~\ref{ch05:rule:master}, Table~\ref{ch05:tab:jcs} flags it as a
deliberate JCS divergence, and the determinism fixture family includes
known-answer encodings whose recorded digests any conforming
implementation must reproduce byte-for-byte. A known-answer test is the
cheapest possible cross-implementation contract, and the terminator is
exactly the kind of detail it exists to catch.
\end{example}

\section{Placement in the evidence pipeline}\label{ch05:pipeline}

Before the testing discipline, this section walks the encoder's actual
consumers, because the profile's rules were each motivated by a threat
in the abstract and each lands somewhere concrete in the kernel.

\emph{Semantic references} (Chapter 3) carry a definition digest: the
hash of the exact defining text of the referenced concept in its
identified source. The dual-identity discipline of
Section~\ref{ch05:byteid} governs how that source text is identified
across platforms, and the identifier profile governs the reference's
namespace and identifier components. A reference is thus a consumer of
this chapter twice over---its target is named by a digest, and its own
rendered form obeys the identifier tier.

\emph{Boundary carriers} (Chapter 4) are encoded through their declared
view, as worked in Section~\ref{ch05:worked}; their digests appear
inside every artifact that embeds a boundary, which is every artifact.
The tagging rules guarantee that a boundary's digested content cannot
be counterfeited from look-alike string data, and the single-route rule
guarantees that two components digesting the same boundary agree.

\emph{Registry records} (Chapter 6) are the heaviest consumers: every
record carries the canonical digest of its payload, recomputed and
compared during the reconstruction rule at every load, resolution, and
export. The registry layer additionally enforces the parse-side
counterparts of this chapter's rules---duplicate-key rejection, text
admissibility of ledger strings, reserved-tag rejection in ledger
data---so that content entering by parse meets the same bar as content
entering by construction.

\emph{Compatibility receipts and replay} (Chapters 7 through 9) reduce
to digest comparison end to end: compiler bytes are hashed before
execution, raw and LF compiler identities are recorded as separate
fields per Rule~\ref{ch05:rule:dual}, result payloads are canonically
digested, and the replay requirement compares fresh digests against
recorded ones before any structural decision passes. Here the encoding
carries its largest responsibility: a nondeterministic encoder would
make honest replays fail and---far worse---a punning encoder could make
a dishonest replay pass. The entire force of the replay gate rests on
Proposition~\ref{ch05:prop:injective}.

\begin{threat}[Profile drift across versions]\label{ch05:threat:drift}
A future revision of the encoding profile---a changed limit, a new
tagged type, a different key ordering---silently invalidates every
stored digest computed under the old profile, or, worse, coexists with
it so that the same value has two ``canonical'' digests in circulation.
\emph{Answered by} treating the profile itself as an identified
artifact: the profile's rules and limits are part of the kernel's
specified surface (Remark~\ref{ch05:rem:limits}), receipts record the
profile identity alongside their digests, and any profile change is a
version change that new receipts name explicitly. Old digests remain
verifiable under the old profile, new digests under the new, and no
digest is ever compared across profiles. The known-answer fixtures of
Section~\ref{ch05:testing} are the tripwire: a profile change that was
not intended as one fails the recorded encodings immediately.
\end{threat}

The pipeline view also clarifies a division of responsibility that
recurs in review discussions. The encoder guarantees that equal values
produce equal bytes; it does not guarantee that the \emph{right} values
were encoded. A receipt that digests the wrong compiler bytes, or a
record that faithfully encodes a false statement, passes every rule in
this chapter. Those failures belong to the layers that choose what to
encode---the protocol disciplines of Chapter 7 and the registry
contracts of Chapter 6---and the encoding's contribution is to make
their choices exact, durable, and comparable, so that a wrong choice is
at least a definite, citable, refutable one.

\section{Testing discipline}\label{ch05:testing}

The encoding profile is only as credible as the fixtures that pin it
down, and the bound Gate-4 suite organizes them into families that
mirror the rule set. The families are recorded explicitly here,
both because they are the concrete referent of every ``blocked by''
clause above and because their structure is itself a design artifact:
each family tests one clause in isolation, so a future regression names
its rule.

\begin{checkclass}[Encoding fixture families]\label{ch05:check:fixtures}
(i) \emph{Acceptance and determinism}: admissible payloads encoded
repeatedly and across permuted construction orders, asserting
byte-identical output, plus known-answer encodings with recorded
digests. (ii) \emph{Rejection, per clause}: sets and frozen sets;
floats; NaN; infinities (including non-finite decimals); non-string
keys; the reserved tag key in data positions; control, format, and
surrogate code points in text; oversized integers. (iii) \emph{Tagging}:
the decimal/string/integer triple of Example~\ref{ch05:ex:decimal};
precision-preserving decimals; tuple-versus-list; enumeration-versus-
value; boolean-versus-integer. (iv) \emph{Bounds and cycles}: payloads
at and just beyond the depth, node, and byte limits; self-referential
and mutually referential structures asserting the typed cycle error;
shared acyclic substructure asserting by-value encoding. (v) \emph{Dual
identity}: artifacts in CRLF, LF, CR, mixed, and no-line-ending forms,
asserting raw and comparison digests, census counts, the census
arithmetic of Proposition~\ref{ch05:prop:census}, and profile
classification. (vi) \emph{Identifier profiles}: admissible and
inadmissible identifiers across the ASCII, NFKC-fixed-point, and
alphabet clauses, and digest-case normalization. The Gate-4 receipt
records these families within the 61-method focused suite; per the
evidence bounds of Chapter 12, they establish tested behavior on the
fixtures, not conformance over all inputs.
\end{checkclass}

Two testing-discipline points generalize beyond this chapter. First,
rejection tests are first-class citizens, equal in number and care to
acceptance tests: a fail-closed system's contract \emph{is} mostly its
rejections, and a suite that only exercises the happy path verifies the
least important half of the design. Second, determinism is tested by
adversarial construction rather than by repetition alone---the same
logical payload built through different insertion orders, merge
histories, and container provenances---because repetition under one
construction path can pass forever while a second path diverges. Both
points are inherited by the registry fixtures of Chapter 6, which
layer parse-level rejections on top of this chapter's encoding-level
ones.

\section{Summary and boundaries}\label{ch05:summary}

The canonical evidence encoding is the kernel's answer to a simple
observation with unforgiving consequences: every comparison the kernel
ever makes is a comparison of bytes, so the value-to-bytes function must
be engineered to the same standard as any other trust-bearing component.
The v0.1 profile achieves this by restriction and explicitness---sorted
deterministic JSON as the surface; a rejection set (unordered containers,
floats, NaN, infinities, dirty text, oversize structure) that converts
every source of nondeterminism or punning into a typed error; a tagged
extension that keeps exact decimals, tuples, and enumerations distinct
where JSON would conflate them; a single declared route per public type;
and a dual raw/comparison identity with a line-ending census for the
artifacts the kernel reads rather than writes. Where the profile follows
RFC 8785 it does so knowingly, and where it departs---numbers above all
---the departure is a stated consequence of the evidence threat model
rather than an accident of implementation.

\begin{boundary}[What the encoding does not establish]\label{ch05:bd:close}
Determinism and injectivity on the admitted domain are properties of the
encoding function, evidenced by the bounded Gate-4 fixtures; they are not
proofs over all inputs, not a cryptographic analysis of SHA-256, and not
statements about layers above or below the encoder. A canonical digest
identifies bytes. It does not validate the semantics of what the bytes
say, does not establish who produced them, and---per
Invariant~\ref{ch04:inv:selfdigest}---never functions as authorization.
Those limits are not defects of the profile; they are its honest
perimeter, and the registry and decision layers of Chapters 6 and 9 are
built against exactly that perimeter.
\end{boundary}
